%!TEX program = lualatex
\documentclass[ngerman,headheight=70pt]{scrartcl}
\usepackage[ngerman]{babel}
\usepackage{blindtext}
\usepackage[utf8]{luainputenc}
\usepackage{scrlayer-scrpage}
\usepackage[T1]{fontenc}
\usepackage{fontspec}
\usepackage{multicol}
\usepackage{enumitem}
\usepackage{epigraph}
\usepackage[german=quotes]{csquotes}
\usepackage[bottom=2cm,footskip=8mm]{geometry}
\setmainfont{Source Sans Pro}
\MakeOuterQuote{"}

\setlength{\parindent}{0pt}
\setlength{\parskip}{6pt}

\pagestyle{plain.scrheadings}
\thispagestyle{scrheadings}
\lohead{\vspace*{-1.5cm}\includegraphics[scale=0.4]{../images/stupa_logo_rot.pdf}}
\rohead{\\Präsidium des\\Studierendenparlaments}

\lofoot[Universität Hamburg · Präsidium des Studierendenparlaments\\\hrule Von-Melle-Park 5 · D-20146 Hamburg]{Universität Hamburg · Präsidium des Studierendenparlaments\\\hrule Von-Melle-Park 5 · D-20146 Hamburg}
\cofoot[]{}
\cohead[\thepage]{}

\renewcommand{\thesubsection}{TOP \arabic{subsection}}
\renewcommand{\thesubsubsection}{\arabic{subsubsection}.}
\setcounter{subsection}{-1}

\begin{document}
    UHH · StuPa-Präsidium · Von-Melle-Park 5 · D-20146 Hamburg

    \section*{Beschluss des Studierendenparlamentes vom 10. November 2016}

    Das StuPa begrüßt den Beschlus des Akademischen Senats vom 13. Oktober
    2016 zum Aufruf zur Mahnwache zum 78. Jahrestag der Pogromnacht und das
    Bemühen, daraus einen gemeinsamen Aufruf von StuPa und AS zu machen. Das
    StuPa ist bestrebt, dass zur Mahnwache im nächsten Jahr ein gemeinsamer
    Aufruf gelingt.

    \subsection*{Erinnern, Lernen, Gestalten}
    \subsubsection*{Aufruf zur Mahnwache zum 78. Jahrestag der Pogromnacht}

    \blockquote[Gotthold Ephraim Lessing: Nathan der Weise, 1779.]{
        Tempelherr: "Wenn aber nun das Kind, \\
        Erbarmte seiner sich der Jude nicht, \\
        Vielleicht im Elend umgekommen wäre?" \\
        Patriarch: "Tut nichts! der Jude wird verbrannt!"\\
    }

    \blockquote[Zygmunt Bauman (Prof. Em. University of Leeds): Die Welt in Panik. Wie die Angst vor Migranten geschürt wird, in: Blätter für deutsche und internationale Politik, 10/16, S. 49.]{
        Die Menschheit befindet sich in der Krise –
        und es gibt keinen anderen Ausweg aus dieser Krise als die
        Solidarität zwischen den Menschen.
    }

    Die Universität Hamburg hat ihr geschichtliches Umfeld im ehemaligen Zentrum
    jüdischen Lebens: im Grindelviertel. Seit dem Ende des 19. Jahrhunderts
    entfaltete sich hier die jüdische Gemeinde und erreichte in den
    republikanischen Weimarer Jahren eine Blüte. Ihr Zentrum bildete die
    "Bornplatz-Synagoge" (auf dem heutigen Joseph- Carlebach-Platz) mit der
    benachbarten Talmud-Tora-Schule im Grindelhof. Das Grindelviertel war ein
    belebter, kulturell pluraler Stadtteil, Ort des urbanen Lebens und der
    Integration, der Bildung und Aufklärung. Die Universität, 1919 gegründet
    und zunächst einer demokratisch-humanistischen Bildungstradition zugewandt,
    war ein Teil dieses Lebens. Doch schon vor der Machtübertragung an die Nazis
    1933 mehrten sich antisemitische Vorfälle, auch an der Universität. Widerspruch
    und Widerstand gegen die rassistische Menschenhatz wurden in den Untergrund
    gedrängt. So erhob sich kein Protest mehr, als am 9.11.1938 in unmittelbarer
    Nachbarschaft die Synagoge geschändet wurde und auch nicht, als neben dem
    Hauptgebäude, in der Moorweidenstraße, ab Herbst 1941 die als Juden verfolgten
    Bürgerinnen und Bürger im enteigneten "Logenhaus" für die Deportation in die
    Vernichtungslager gesammelt wurden.

    Der Antisemitismus der Nazis diente dazu, in der tiefen Gesellschaftskrise
    seit 1928 der Bevölkerung Sündenböcke darzubieten, die tiefen sozialen
    Verwerfungen so zu verfestigen und mit Einschüchterung, Terror und völkischer
    Ideologie einen Raubkrieg vorzubereiten, der 60 Millionen Menschenleben
    forderte. Demütigung, Vertreibung und Ignoranz gegen rassisch verfolgte,
    gegen bürgerlich-humanistische, sozialistische und kommunistische
    Kommilitoninnen und Kommilitonen, Kolleginnen und Kollegen, Nachbarinnen
    und Nachbarn prägten das düsterste Kapitel der Universitätsgeschichte. An
    ihre so geschmähten Mitglieder erinnert die Universität aus gutem Grund:
    nicht nur das ihnen zugefügte Leid abzulehnen, sondern auch ihr engagiertes
    Wirken für eine menschenwürdige Gesellschaft aufzugreifen, ist für heute
    bedeutsam.

    Nie wieder sollen Bildung und Wissenschaft aus Neid, Konkurrenz und
    Vorurteilen, aus Gleichgültigkeit, Ressentiment und Opportunitätsdenken
    mitverantwortlich werden an Ausgrenzung, Verfolgung, Mord und Krieg. Frieden,
    Gerechtigkeit und Humanität sollen Reflexion, Diskurs und gesellschaftliche
    Praxis der Universität nachhaltig orientieren. Weltoffenheit, Inklusion und
    Demokratie sollen ihre Kultur prägen.

    Wir erinnern, für ein besseres Leben.

    {\centering

    \textbf{Mahnwache}

    \textbf{anlässlich des 78. Jahrestags der Reichspogromnacht \\
    am Mittwoch, den 9. November 2016 von 15.30 – 17.00 Uhr \\
    Joseph-Carlebach-Platz (Grindelhof).}\\}
    \vspace{0.3cm}

    Es sprechen: Peggy Parnass, Michael Heimann (Jüdischen Gemeinde Hamburg),
    Wolfgang Seibert (Jüdische Gemeinde Pinneberg – angefragt), Norma van der
    Walde (VVN-BdA), Prof. Dr. Susanne Rupp (Universität Hamburg)

    Moderation: Traute Springer-Yakar - VVN-BdA Hamburg

    Veranstalter: Vereinigung der Verfolgten des Naziregimes – Bund der
    Antifaschisten (VVN-BdA), Jüdische Gemeinde Hamburg, Universität Hamburg

\end{document}
